% ==============================================================================
% File: lf_paper.tex
% Bài báo: Thinking with Looped Flows
% Tác giả: Ayhan Suleymanzade, Chanhyuk Lee, Floor Eijkelboom, Nicholas M. Boffi,
%          İsmail İlkan Ceylan, Jinwoo Kim (arXiv:2609.11801, 2026)
% Nguồn số liệu: reports/2609.11801_looped_flows.md
% ==============================================================================

\section{Dòng vòng lặp Looped Flows}

% ------------------------------------------------------------------------------
% Slide 0: Tổng quan bài báo
% ------------------------------------------------------------------------------
\begin{frame}{Looped Flows: Tổng quan bài báo}
    \begin{itemize}
        \item \textbf{Tên bài báo:} \textit{Thinking with Looped Flows}
        \item \textbf{Tác giả và tổ chức:} Ayhan Suleymanzade, Chanhyuk Lee, Jinwoo Kim, Nicholas M. Boffi et al., EPFL, KAIST, CMU và các nơi khác, arXiv:2609.11801
        \item \textbf{Ý tưởng cốt lõi:} Huấn luyện mô hình vòng lặp bằng một chuỗi mục tiêu khử nhiễu cục bộ, thay vì lan truyền ngược qua nhiều bước lặp.
        \item \textbf{Đóng góp nổi bật:}
        \begin{itemize}
            \item \textbf{Bộ khử nhiễu có trạng thái:} vừa khử nhiễu, vừa cập nhật trạng thái lặp. Nhiễu giảm dần và nhiễu dùng chung buộc trạng thái lặp mang thông tin hữu ích cho bước sau.
            \item \textbf{Suy luận bằng tích phân dòng xác suất:} tăng tính toán qua lưới thời gian mịn hơn, và cho nhiều lời giải hợp lệ từ các mẫu nhiễu khác nhau.
            \item \textbf{Kết quả:} ARC-AGI-1 đạt 58,8\% (TRM 44,6\%), ARC-AGI-2 đạt 12,2\% (TRM 7,8\%), Sudoku-Extreme đạt 97,9\%, chỉ với 5 đến 7M tham số.
        \end{itemize}
    \end{itemize}
\end{frame}

% ------------------------------------------------------------------------------
% Slide 1: Bối cảnh và thách thức
% ------------------------------------------------------------------------------
\begin{frame}{1. Looped Flows: Bối cảnh và thách thức}
    \textbf{Bối cảnh:} Bài toán khó cần nhiều tính toán hơn lúc suy luận. Mô hình vòng lặp cập nhật trạng thái ẩn nhiều lần để tích lũy tính toán.
    \vspace{0.1cm}

    \textbf{Ba khoảng trống của các hướng hiện có:}
    \begin{itemize}
        \item \textbf{Mô hình vòng lặp như HRM, TRM, FPRM:} lan truyền ngược qua toàn bộ chuỗi lặp vừa tốn bộ nhớ vừa bất ổn, nên phải chặn gradient. Mất mát ở bước $\ge i$ không gửi được tín hiệu cho trạng thái $z_{<i}$.
        \item \textbf{Mô hình dòng và khuếch tán:} học khử nhiễu cục bộ, không cần lan truyền ngược qua quỹ đạo, nhưng không có trạng thái lặp, nên yếu với tính toán tuần tự.
        \item \textbf{Mô hình vòng lặp ngẫu nhiên như PTRM, EqR, GRAM:} thêm tính ngẫu nhiên nhưng vẫn dựa vào gradient cục bộ để xây trạng thái lặp hữu ích toàn cục.
    \end{itemize}
    \vspace{0.1cm}
    \textbf{Hậu quả:} chuỗi lặp bất ổn, rơi vào điểm hút giả hoặc không hội tụ, ngay cả ở bài toán đơn giản. Hướng giải quyết: kết hợp mục tiêu khử nhiễu cục bộ với trạng thái lặp.
\end{frame}

% ------------------------------------------------------------------------------
% Slide 2: Kiến thức nền
% ------------------------------------------------------------------------------
\begin{frame}{2. Looped Flows: Kiến thức nền}
    \textbf{Mô hình vòng lặp:} $z_i = f(z_{i-1}; c)$, đầu đọc $\hat{x}_i = g(z_i)$. Thực tế huấn luyện với chặn gradient $\mathrm{sg}$:
    $$\mathcal{L}_{\text{loop}} = \mathbb{E}_{c,x}\Big[\sum_i \mathrm{CE}(\hat{x}_i, x)\Big], \qquad z_i = f\big(\mathrm{sg}(z_{i-1}); c\big)$$
    \begin{itemize}
        \item Điểm yếu: mất mát ở bước $\ge i$ không truyền được về $z_{<i}$. Quá trình lặp xác định chỉ cho một lời giải.
    \end{itemize}
    \vspace{0.1cm}
    \textbf{Khớp dòng cho dữ liệu phân loại:} nội suy $I_t = (1-t)x_0 + t x_1$, với $x_0 \sim \mathcal{N}(0,\sigma^2 I)$.
    \begin{itemize}
        \item Bộ khử nhiễu $D_t(x;c) = \mathbb{E}[x_1 \mid I_t = x, c]$, vận tốc $b_t(x;c) = \frac{D_t(x;c) - x}{1-t}$.
        \item Hàm mục tiêu $\mathcal{L}_{\text{flow}} = \mathbb{E}\,[\mathrm{CE}(\hat{x}_t, x_1)]$ với $t \sim \mathcal{U}[0,1]$: mất mát tại $t$ không phụ thuộc tính toán ở $t' < t$, nên không cần lan truyền ngược qua quỹ đạo.
        \item Điểm yếu: không có trạng thái, khó học tính toán tuần tự.
    \end{itemize}
\end{frame}

% ------------------------------------------------------------------------------
% Slide 3: Bộ khử nhiễu có trạng thái
% ------------------------------------------------------------------------------
\begin{frame}{3. Looped Flows: Bộ khử nhiễu có trạng thái}
    \textbf{Ý tưởng:} Biến bộ khử nhiễu thành có trạng thái, vừa khử nhiễu trạng thái dòng $x_{t_i}$, vừa cập nhật trạng thái lặp $z_{t_i}$:
    $$\big(\hat{x}_{t_i},\, z_{t_{i+1}}\big) = \hat{D}_{t_i}\big(x_{t_i},\, z_{t_i};\, c\big)$$
    \begin{itemize}
        \item Suy luận dùng bước Euler tiến trên lưới thời gian $0 = t_0 < \dots < t_n = 1$:
        $$x_{t_{i+1}} = x_{t_i} + (t_{i+1} - t_i)\,\frac{\hat{x}_{t_i} - x_{t_i}}{1 - t_i}$$
        \item \textbf{Kiến trúc:} xây trên TRM, thêm lớp chiếu cho nội suy nhiễu và nhúng thời gian, thêm khoảng 5\% tham số. Trạng thái lặp gồm hai thành phần, đặc trưng ẩn $h$ để giải mã dự đoán và thành phần $\ell$ để cập nhật $h$. Chỉ lan truyền ngược qua chu kỳ cuối của mỗi lần gọi.
        \item \textbf{Quy mô:} 5M tham số cho Sudoku với MLP-Mixer, 7M tham số cho các tác vụ còn lại với Transformer.
    \end{itemize}
\end{frame}

% ------------------------------------------------------------------------------
% Slide 4: Huấn luyện
% ------------------------------------------------------------------------------
\begin{frame}{4. Looped Flows: Căn chỉnh theo thời gian}
    \textbf{Hàm mục tiêu:} lấy $k+1$ mốc thời gian đều, sắp tăng dần $t_0 < \dots < t_k$, rồi cộng cross-entropy qua các bước:
    $$\mathcal{L}_{\text{LF}} = \mathbb{E}\Big[\sum_{i=0}^{k-1} \mathrm{CE}(\hat{x}_{t_i}, x_1)\Big], \quad (\hat{x}_{t_i}, z_{t_{i+1}}) = \hat{D}_{t_i}\big(I_{t_i}, \mathrm{sg}(z_{t_i}); c\big)$$
    \textbf{Hai thiết kế then chốt:}
    \begin{itemize}
        \item \textbf{Nhiễu giảm dần:} mức nhiễu $1 - t_i$ giảm theo thời gian, thành chuỗi bài khử nhiễu từ khó đến dễ.
        \item \textbf{Nhiễu dùng chung:} cùng một cặp $(x_0, x_1)$ cho mọi $t_i$ trong một lần chạy. Các bước kề nhau có đầu vào gần nhau, nên trạng thái lặp $z_{t_i}$ vẫn hữu ích ở bước sau dù gradient chỉ phủ một bước.
    \end{itemize}
    \textbf{Mục tiêu giả:} khi dữ liệu ít, ví dụ 1000 mẫu Sudoku, nội suy ở $t$ lớn gần như lộ $x_1$ và gây overfit. Thay $x_1$ trong nội suy bằng dự đoán trước đó $\mathrm{sg}(\hat{x}_{t_{i-1}})$, xác suất dùng tăng tuyến tính từ 0 đến 1 trong 20 nghìn bước. Mất mát vẫn dùng $x_1$ thật.
\end{frame}

% ------------------------------------------------------------------------------
% Slide 5: ACT và chi tiết huấn luyện
% ------------------------------------------------------------------------------
\begin{frame}{5. Looped Flows: Thời gian tính toán thích ứng và cấu hình}
    \begin{itemize}
        \item \textbf{Vấn đề:} nhiễu giảm dần làm mất mát khử nhiễu giảm đơn điệu và bão hòa, các bước cuối không còn tín hiệu học.
        \item \textbf{Đầu dừng $q$} là bộ phân loại nhị phân xem dự đoán hiện tại đã đúng chưa:
        $$\mathcal{L}_{\text{ACT}} = \mathbb{E}\Big[\sum_{i} \lambda\,\mathrm{BCE}\big(q(z_{t_{i+1}}),\, \mathbb{1}(x_1 = \mathrm{round}(\hat{x}_{t_i}))\big)\Big]$$
        Khi $q > 1/2$, các bước sau bị bỏ qua lúc huấn luyện. Tổng mất mát $= \mathcal{L}_{\text{LF}} + \mathcal{L}_{\text{ACT}}$, $\lambda = 0{,}5$.
        \item \textbf{Cấu hình:} $k = 16$ bước, Adam-atan2, tốc độ học đỉnh $10^{-4}$, batch 768, EMA 0,999, cắt gradient 1,0.
        \item \textbf{Tài nguyên:} khoảng 1 đến 5 giờ cho Sudoku, Maze, N-Queens, Graph Coloring. Khoảng 1 đến 2 ngày cho ARC.
    \end{itemize}
\end{frame}

% ------------------------------------------------------------------------------
% Slide 6: Suy luận
% ------------------------------------------------------------------------------
\begin{frame}{6. Looped Flows: Cơ chế suy luận}
    \begin{itemize}
        \item \textbf{Tích phân xác định bằng ODE:} bước Euler tiến, mỗi lần gọi bộ khử nhiễu cập nhật cả $x$ và $z$.
        \item \textbf{Tích phân ngẫu nhiên bằng SDE:} tương đương SDE của dòng xác suất, thực hiện bằng lùi nhiễu. Mỗi bước đẩy $x_{t_i}$ lùi về thời điểm $s = a\,t_i$, thêm nhiễu mới, rồi bước Euler tới $t_{i+1}$. Hệ số $\gamma = 0$ quay về ODE, $\gamma > 0$ thêm tính ngẫu nhiên.
        \item \textbf{Mở rộng tính toán lúc suy luận:}
        \begin{itemize}
            \item \textbf{Lưới thời gian mịn hơn lúc huấn luyện.} Sudoku huấn luyện 16 bước nhưng suy luận 128 bước.
            \item \textbf{Tổ hợp chọn mẫu theo điểm Q cao nhất:} chạy nhiều quỹ đạo độc lập, thường 5, chọn dự đoán có xác suất thành công cao nhất theo đầu dừng $q$.
        \end{itemize}
        \item \textbf{Nhiều lời giải hợp lệ:} xuất phát từ các mẫu nhiễu $x_0$ khác nhau, vận chuyển xác suất tới các mode khác nhau của $p_{\text{data}}(x \mid c)$. Mô hình vòng lặp xác định chỉ cho một lời giải mỗi bài.
    \end{itemize}
\end{frame}

% ------------------------------------------------------------------------------
% Slide 7: Kết quả chính
% ------------------------------------------------------------------------------
\begin{frame}{7. Looped Flows: Kết quả trên các bài toán suy luận}
    \textbf{Độ chính xác lời giải, phần trăm. ARC dùng pass@2. Looped Flows chạy ba hạt giống.}
    \begin{table}[h]
        \centering
        \small
        \begin{tabular}{lcccc}
            \toprule
            \textbf{Phương pháp} & \textbf{Sudoku-Extreme} & \textbf{Maze-Hard} & \textbf{ARC-AGI-1} & \textbf{ARC-AGI-2} \\
            \midrule
            HRM, 27M & 55,0 & 74,5 & 40,3 & 5,0 \\
            TRM, 5 đến 7M & 87,4 & 85,3 & 44,6 & 7,8 \\
            FPRM, 7M & 94,2 & \textbf{87,0} & 47,5 & 6,2 \\
            GRAM, 10M & -- & -- & 52,0 & 11,1 \\
            \midrule
            \textbf{Looped Flows, 5 đến 7M} & \textbf{97,9 $\pm$ 0,4} & 86,7 $\pm$ 1,1 & \textbf{58,8 $\pm$ 1,8} & \textbf{12,2 $\pm$ 1,9} \\
            \bottomrule
        \end{tabular}
    \end{table}
    \vspace{-0.2cm}
    \begin{small}
    \begin{itemize}
        \item Tốt nhất ở 3 trong 4 bài; Maze-Hard chỉ ngang FPRM, bài báo không giải thích.
        \item Tổ hợp 5 quỹ đạo: Sudoku 99,3, Maze 86,9, ARC-AGI-1 59,5, ARC-AGI-2 12,2. Cạnh tranh với PTRM dùng 100 quỹ đạo và EqR dùng 128 quỹ đạo.
        \item Độ lệch chuẩn trên ARC khoảng 1,8 đến 1,9 điểm: kết quả nhạy với hạt giống.
    \end{itemize}
    \end{small}
\end{frame}

% ------------------------------------------------------------------------------
% Slide 8: Mở rộng lúc suy luận và độ ổn định
% ------------------------------------------------------------------------------
\begin{frame}{8. Looped Flows: Mở rộng lúc suy luận và độ ổn định}
    \textbf{Tăng số bước tích phân trên Sudoku:}
    \begin{center}
        \small
        \begin{tabular}{lccccc}
            \toprule
            \textbf{Số bước} & 8 & 16 & 32 & 64 & 128 \\
            \midrule
            \textbf{Độ chính xác (\%)} & 74,5 & 85,2 & 93,0 & 96,7 & 97,9 \\
            \bottomrule
        \end{tabular}
    \end{center}
    \begin{itemize}
        \item Vượt GRAM tại 32 bước.
    \end{itemize}
    \vspace{0.1cm}
    \textbf{Độ ổn định của chuỗi lặp, trên khoảng 65 nghìn bài Sudoku-Extreme:}
    \begin{itemize}
        \item TRM thất bại ở 12,6\% số bài: 88,3\% do không hội tụ, 11,7\% do điểm hút giả.
        \item Looped Flows cứu được \textbf{90,9\%} các bài TRM thất bại: 89,9\% trường hợp không hội tụ và 98,0\% trường hợp điểm hút giả.
    \end{itemize}
\end{frame}

% ------------------------------------------------------------------------------
% Slide 9: Đa lời giải
% ------------------------------------------------------------------------------
\begin{frame}{9. Looped Flows: Bài toán đa lời giải}
    \textbf{So sánh với GRAM. Độ phủ là tỉ lệ lời giải thật được tìm thấy.}
    \begin{table}[h]
        \centering
        \small
        \begin{tabular}{llcc}
            \toprule
            \textbf{Bài toán} & \textbf{Chỉ số} & \textbf{GRAM} & \textbf{Looped Flows} \\
            \midrule
            N-Queens 8$\times$8 & Độ chính xác & 99,7 & \textbf{99,9} \\
            & Độ phủ & 90,3 & \textbf{91,4} \\
            N-Queens 10$\times$10 & Độ chính xác & 96,1 & 94,4 \\
            & Độ phủ & \textbf{74,3} & 61,5 \\
            \midrule
            Tô màu đồ thị 8 đỉnh & Số xung đột $\downarrow$ & 2,7 & \textbf{0,7} \\
            & Độ phủ & 85,8 & \textbf{89,4} \\
            Tô màu đồ thị 10 đỉnh & Số xung đột $\downarrow$ & 3,3 & \textbf{1,0} \\
            & Độ phủ & 51,3 & \textbf{55,2} \\
            \bottomrule
        \end{tabular}
    \end{table}
    \vspace{-0.2cm}
    \begin{itemize}
        \item Trên N-Queens 10$\times$10, GRAM vẫn hơn ở cả độ chính xác và độ phủ. Looped Flows tốt nhất ở 6 trong 8 chỉ số.
    \end{itemize}
\end{frame}

% ------------------------------------------------------------------------------
% Slide 10: Nghiên cứu loại bỏ thành phần
% ------------------------------------------------------------------------------
\begin{frame}{10. Looped Flows: Nghiên cứu loại bỏ thành phần}
    \textbf{Khi huấn luyện, trên ARC-AGI-1 và ARC-AGI-2:}
    \begin{center}
        \small
        \begin{tabular}{lcc}
            \toprule
            \textbf{Cấu hình} & \textbf{ARC-AGI-1} & \textbf{ARC-AGI-2} \\
            \midrule
            \textbf{Đầy đủ} & \textbf{58,8} & \textbf{12,2} \\
            Bỏ điều kiện thời gian & 56,4 & 9,9 \\
            Bỏ nội suy & 51,5 & 9,9 \\
            Bỏ cả hai & 43,6 & 5,0 \\
            Bỏ nhiễu giảm dần & 51,6 & 9,9 \\
            Bỏ nhiễu dùng chung & 56,4 & 10,8 \\
            \bottomrule
        \end{tabular}
    \end{center}
    \begin{small}
    \begin{itemize}
        \item Nhiễu giảm dần quan trọng nhất; nhiễu dùng chung ít hơn nhưng vẫn giúp. Bỏ cả hai thành phần dòng thì ngang TRM (43,6 so với 44,6).
        \item \textbf{Lúc suy luận:} SDE với $\gamma = 5$ nhất quán tốt hơn ODE, ví dụ Sudoku 97,9 so với 97,6, ARC-AGI-1 58,8 so với 57,5.
        \item \textbf{Trên Sudoku:} dòng không vòng lặp ghi nhớ dữ liệu, tổng quát hóa kém; tự điều kiện hóa tốt hơn; Looped Flows tốt nhất.
    \end{itemize}
    \end{small}
\end{frame}

% ------------------------------------------------------------------------------
% Slide 11: Hạn chế, phản biện và kết luận
% ------------------------------------------------------------------------------
\begin{frame}{11. Looped Flows: Hạn chế, phản biện và kết luận}
    \small
    \textbf{Hạn chế và điểm cần thận trọng:}
    \begin{itemize}
        \item Chưa có lý thuyết cho căn chỉnh theo thời gian. Nhiễu dùng chung về lý thuyết cho phép đường tắt triệt nhiễu; bài báo chỉ cho thấy gián tiếp rằng mô hình không dùng nó.
        \item Lệch phân phối giữa huấn luyện dùng nội suy $I_{t_i}$ và suy luận dùng $x_{t_i}$ từ bước Euler trước. Bài báo không bàn tới.
        \item Mô hình 5 đến 10M tham số, chỉ trên câu đố suy luận, nhiều siêu tham số riêng cho từng tác vụ. Nghiên cứu loại bỏ chỉ trên ARC.
        \item Huấn luyện vẫn cần chạy tuần tự $k$ bước, chưa phải huấn luyện không mô phỏng.
    \end{itemize}
    \textbf{Kết luận:}
    \begin{itemize}
        \item Mục tiêu khử nhiễu cục bộ cộng căn chỉnh theo thời gian huấn luyện được trạng thái lặp ổn định, kể cả khi gradient chỉ phủ một bước.
        \item Cho phép mở rộng tính toán lúc suy luận và sinh nhiều lời giải. \textbf{Hướng tiếp theo:} huấn luyện không mô phỏng, quy mô lớn hơn.
    \end{itemize}
\end{frame}
